% Part: history
% Chapter: biographies
% Section: ernst-zermelo

\documentclass[../../../include/open-logic-section]{subfiles}

\begin{document}

\olfileid{his}{bio}{zer}

\olsection{एर्न्स्ट झर्मेलो}

\olphoto{zermelo-ernst}{एर्न्स्ट झर्मेलो}

एर्न्स्ट झर्मेलो यांचा जन्म २७ जुलै १८७१ रोजी जर्मनीतील बर्लिन येथे झाला. त्यांना पाच बहिणी होत्या; मात्र कुटुंबातील सदस्यांची तब्येत नाजूक असल्याने त्यांपैकी केवळ तीनच जणी प्रौढ वयापर्यंत जगल्या. झर्मेलो लहान असतानाच त्यांच्या आईवडिलांचेही निधन झाले आणि वयाच्या सतराव्या वर्षी ते व त्यांची भावंडे अनाथ झाली. झर्मेलो यांना कलांत, विशेषतः कवितेत, मोठा रस होता. ते तीक्ष्ण बुद्धीचे, चतुर आणि टीकात्मक स्वभावाचे म्हणून ओळखले जात. त्यांच्या सर्वाधिक प्रसिद्ध गणितीय कामगिरीत १९०४ मध्ये निवडीचे स्वयंसिद्धक मांडणे आणि १९०८ मध्ये संच उपपत्तीची स्वयंसिद्धकीय मांडणी करणे यांचा समावेश होतो.

विद्यापीठात झर्मेलो यांच्या आवडी विविध होत्या. त्यांनी भौतिकशास्त्र, गणित आणि तत्त्वज्ञान यांचे अभ्यासक्रम घेतले. हेरमान श्वार्त्स यांच्या मार्गदर्शनाखाली त्यांनी १८९४ मध्ये बर्लिन विद्यापीठात \emph{Investigations in
  the Calculus of Variations} हा प्रबंध पूर्ण केला. १८९७ मध्ये पुढील शिक्षणासाठी त्यांनी गॉटिंगेन विद्यापीठात जाण्याचे ठरवले; तेथे डेव्हिड हिल्बर्ट यांच्या गणिताच्या पायांवरील कार्याचा त्यांच्यावर मोठा प्रभाव पडला. १८९९ मध्ये ते प्राध्यापकपदासाठी पात्र झाले, पण प्रत्यक्ष पद अकरा वर्षांनीच मिळाले. त्यांच्या विचित्र वागण्यामुळे आणि “अस्वस्थ उतावळेपणा”मुळे हा विलंब झाला असावा, असा स्रोताचा अंदाज आहे.

अखेरीस १९१० मध्ये झर्मेलो यांना झुरिक विद्यापीठात वेतन मिळणारे प्राध्यापकपद मिळाले; मात्र क्षयरोगामुळे १९१६ मध्ये त्यांना निवृत्त व्हावे लागले. बरे झाल्यावर १९२१ मध्ये त्यांना फ्रायबुर्ग विद्यापीठात मानद प्राध्यापकपद देण्यात आले. या काळात त्यांनी गणिताच्या पायांवर काम केले. थोराल्फ स्कोलेम आणि कुर्ट गोडेल यांच्या कार्यावर ते नाराज झाले आणि त्यांनी आपल्या लेखांत त्यांच्या दृष्टिकोनांवर जाहीर टीका केली. लोकप्रियता नसणे आणि जर्मनीत हिटलरच्या सत्तारूढ होण्यास त्यांचा विरोध असणे, या कारणांमुळे १९३५ मध्ये त्यांना फ्रायबुर्गमधील पदावरून काढून टाकण्यात आले, असे स्रोत सांगतो.

झर्मेलो यांच्या आयुष्याचा उत्तरार्ध एकाकीपणात गेला. १९३५ मध्ये पदावरून काढण्यात आल्यानंतर त्यांनी गणिताचे काम सोडले. ते ग्रामीण भागात जाऊन साधेपणाने राहू लागले. १९४४ मध्ये त्यांनी विवाह केला. त्यांची दृष्टी मंदावत असल्याने ते पत्नीवर पूर्णपणे अवलंबून झाले. १९५१ पर्यंत त्यांची दृष्टी पूर्णपणे गेली होती. २१ मे १९५३ रोजी जर्मनीतील ग्युंटरस्टाल येथे त्यांचे निधन झाले.

\begin{reading}
झर्मेलो यांच्या सविस्तर चरित्रासाठी \citet{Ebbinghaus2015} पाहा. त्यांचे १९०४ आणि १९०८ मधील मूलभूत लेख मूळ जर्मन भाषेत \citep{Zermelo1904,Zermelo1908} येथे उपलब्ध आहेत. भौतिकशास्त्रावरील लेखनासह त्यांचे संकलित लेख इंग्रजी भाषांतरात \citep{Ebbinghaus2010,Ebbinghaus2013} येथे उपलब्ध आहेत.
\end{reading}

\end{document}
